\section*{Chapter \ref{ch:g7:sediments-and-fossils} --- Sediments and Fossils}

\begin{solution}{exo:g7:sediments-and-fossils:1}
(a) \omterm{def:g7:sediments-and-fossils:fossil}{mould}; (b) \omterm{def:g7:sediments-and-fossils:fossil}{trace fossil}; (c) \omterm{def:g7:sediments-and-fossils:sediment}{sediment}; (d) \omterm{def:g7:sediments-and-fossils:fossil}{cast}; (e) \omterm{def:g7:sediments-and-fossils:layer}{stratum}; (f)
\omterm{def:g7:sediments-and-fossils:fossil}{fossil}.
\end{solution}

\begin{solution}{exo:g7:sediments-and-fossils:2}
(b) the animal dies; (d) \omterm{def:g3:stones-pebbles-sand:sand}{mud} covers the shell quickly; (c) new layers
pile up and the shell turns to stone; (a) \omterm{def:g5:water-shapes-the-land:erosion}{erosion} wears away the layers
above it.
\end{solution}

\begin{solution}{exo:g7:sediments-and-fossils:3}
(a) A warm, shallow, clear \omterm{def:g2:water-on-earth:ocean}{sea}. (b) Shallow water, at the edge of a
\omterm{def:g1:landscapes-around-us:water}{lake}, a \omterm{def:g1:landscapes-around-us:water}{river} or a \omterm{def:g2:water-on-earth:ocean}{sea}. (c) A swamp forest or a marsh. (d) Land, with
plants growing (and dragonflies flying), near water where the \omterm{def:g3:stones-pebbles-sand:sand}{mud}
settled.
\end{solution}

\begin{solution}{exo:g7:sediments-and-fossils:4}
(a) False: most living things leave no \omterm{def:g7:sediments-and-fossils:fossil}{fossil}; they rot or dissolve
before they are buried. (b) True: it is a \omterm{def:g7:sediments-and-fossils:fossil}{trace fossil}. (c) False: the
largest and heaviest grains settle first, the fine \omterm{def:g3:stones-pebbles-sand:sand}{mud} last. (d) True.
\end{solution}

\begin{solution}{exo:g7:sediments-and-fossils:5}
(a) Sandstone \qty{120}{m}, limestone \qty{180}{m}, shale \qty{100}{m}.
(b) $120 + 180 + 100 = \qty{400}{m}$. (c) $180 \div 400 = 0.45$:
\textbf{45~\%}. (d) Its top has been worn away by \omterm{def:g5:water-shapes-the-land:erosion}{erosion}: what is left
is not the whole layer, and how much is left changes from place to place.
\end{solution}

\begin{solution}{exo:g7:sediments-and-fossils:6}
(a) From the bottom up: \omterm{def:g3:stones-pebbles-sand:pebble}{gravel}, then \omterm{def:g3:stones-pebbles-sand:sand}{sand}, then \omterm{def:g3:stones-pebbles-sand:sand}{mud} on top, under the
clear water. (b) The tiny \omterm{def:g3:stones-pebbles-sand:sand}{mud} grains sink very slowly: they are still
falling after an hour. (c) A second set of layers on top of the first:
\omterm{def:g3:stones-pebbles-sand:sand}{sand}, then \omterm{def:g3:stones-pebbles-sand:sand}{mud} again. (d) The \omterm{def:g3:stones-pebbles-sand:pebble}{gravel} at the very bottom.
\end{solution}

\begin{solution}{exo:g7:sediments-and-fossils:7}
Best: (b), a hard shell, buried quickly under fine \omterm{def:g3:stones-pebbles-sand:sand}{mud} in calm water.
Then (d): the shell is hard and buried, but a \omterm{def:g5:water-shapes-the-land:flood}{flood} can also move and
break it. Least: (a), a soft body with no hard parts, on \omterm{def:g3:stones-pebbles-sand:rock}{rocks} where no
\omterm{def:g3:stones-pebbles-sand:sand}{mud} settles to cover it. (c) is not much better: the bone stays
uncovered and is broken down and worn away.
\end{solution}

\begin{solution}{exo:g7:sediments-and-fossils:8}
(a) A hollow with the shape of the shell printed on its walls: a \omterm{def:g7:sediments-and-fossils:fossil}{mould}.
(b) A \omterm{def:g7:sediments-and-fossils:fossil}{cast}, a stone copy of the shell. (c) No: the shell has dissolved;
the \omterm{def:g7:sediments-and-fossils:fossil}{cast} is made of new \omterm{def:g6:rocks-and-minerals:mineral}{minerals}. (d) Water passes easily between the
large grains of \omterm{def:g3:stones-pebbles-sand:sand}{sand}, so it can reach and dissolve the shell; in fine
\omterm{def:g3:stones-pebbles-sand:sand}{mud} it passes much less easily.
\end{solution}

\begin{solution}{exo:g7:sediments-and-fossils:9}
(a) $1.6 \times 1000 = \qty{1600}{m}$. (b) $1600 \div 80 = 20$ layers.
(c) $160 \div 1600 = 0.1$: \textbf{10~\%}.
\end{solution}

\begin{solution}{exo:g7:sediments-and-fossils:10}
(a) No: a foot leaves a print only in soft ground, wet \omterm{def:g3:stones-pebbles-sand:sand}{sand} or \omterm{def:g3:stones-pebbles-sand:sand}{mud}; the
ground became \omterm{def:g3:stones-pebbles-sand:rock}{rock} later. (b) The next wave or \omterm{def:g5:water-shapes-the-land:flood}{flood} washes them away.
(c) New \omterm{def:g7:sediments-and-fossils:sediment}{sediment} had to cover them and fill them quickly, before they
were erased. (d) A \omterm{def:g7:sediments-and-fossils:fossil}{trace fossil}: it shows what the animal did, not a
part of its body.
\end{solution}

\begin{solution}{exo:g7:sediments-and-fossils:11}
(a) $20 + 5 = \qty{25}{mm}$. (b) $1000 \div 25 = 40$ pairs: the \omterm{def:g6:inside-the-earth:layers}{core}
records \textbf{40 \omterm{def:g5:water-shapes-the-land:flood}{floods}}. (c) The \omterm{def:g3:stones-pebbles-sand:sand}{sand} is heavy: it settles soon after
the \omterm{def:g1:landscapes-around-us:water}{river} enters the \omterm{def:g1:landscapes-around-us:water}{lake}, near the \omterm{def:g5:water-shapes-the-land:source-mouth}{mouth}; only a little of it travels
far. (d) Far out, only the fine \omterm{def:g3:stones-pebbles-sand:sand}{mud} is still in the water: it settles
there slowly, when the water is calm.
\end{solution}

\begin{solution}{exo:g7:sediments-and-fossils:12}
(a) Shells are hard: they last long enough to be buried, and many
animals with shells lived on the \omterm{def:g2:water-on-earth:ocean}{sea} floor, where \omterm{def:g3:stones-pebbles-sand:sand}{mud} settles over them.
(b) Soft bodies rot before they can be buried. (c) No: soft animals
almost never leave \omterm{def:g7:sediments-and-fossils:fossil}{fossils}, so the list from \omterm{def:g7:sediments-and-fossils:fossil}{fossils} is incomplete; it
shows mostly the animals with shells or bones.
\end{solution}

\begin{solution}{pb:g7:sediments-and-fossils:1}
\textbf{1.} Sandstone, limestone, shale, chalk (the first at the bottom
of the column).

\textbf{2.} \qty{120}{m}, \qty{180}{m} and \qty{100}{m}; total
\qty{400}{m}.

\textbf{3.} $180 \div 400 = 0.45$: \textbf{45~\%}.

\textbf{4.} Its top has been worn away: the upper part of the layer is
gone, and nobody can measure what is no longer there.

\textbf{5.} In shallow water at the edge of a \omterm{def:g1:landscapes-around-us:water}{lake}, a \omterm{def:g1:landscapes-around-us:water}{river} or a \omterm{def:g2:water-on-earth:ocean}{sea}:
ripple marks form on \omterm{def:g3:stones-pebbles-sand:sand}{sand} in shallow water, and footprints are left in
soft, wet \omterm{def:g3:stones-pebbles-sand:sand}{sand} or \omterm{def:g3:stones-pebbles-sand:sand}{mud}, often on a shore.

\textbf{6.} In a warm, shallow, clear \omterm{def:g2:water-on-earth:ocean}{sea}: the corals.

\textbf{7.} Calm: fine \omterm{def:g3:stones-pebbles-sand:sand}{mud} settles only where the water is almost still.

\textbf{8.} First the place was a shore, with shallow water and soft \omterm{def:g3:stones-pebbles-sand:sand}{sand}
where animals walked. Then the \omterm{def:g2:water-on-earth:ocean}{sea} covered it: a warm, shallow \omterm{def:g2:water-on-earth:ocean}{sea} where
corals grew and shells piled up. Later the water became calmer, perhaps
deeper, and only fine \omterm{def:g3:stones-pebbles-sand:sand}{mud} settled, with the shells of ammonites, animals
related to squids and octopuses.

\textbf{9.} $180 \times 100 = 18\,000$~cm.

\textbf{10.} $120~\text{m} = 12\,000$~cm, and $12\,000 \times 20 =
240\,000$~years.

\textbf{11.} $1~\text{m} = 100$~cm, and $100 \times 100 = 10\,000$~years.

\textbf{12.} $180 \times 10\,000 = 1\,800\,000$: about
\textbf{1\,800\,000 years}, almost two million years. That is
$1\,800\,000 \div 240\,000 = 7.5$ times longer than the sandstone.
\end{solution}
